\chapter{弦乐配器：把五条声部写成一支队伍}
\label{ch:strings}

\classquote{声相应，故生变；变成方，谓之音。}{《礼记·乐记》}{声音彼此呼应才生出变化，变化有了章法才成其为音乐。五个弦乐声部正是这样互相接话的。弦乐好听，不是因为“弦乐音色高级”，而是因为五个声部各有各的活：谁唱歌、谁垫和声、谁守低音、谁在中间填空。}

\section{弦乐组先是一支队伍，不是一件乐器}
\label{sec:string-section}

\term{弦乐组}{string-section}由第一小提琴、第二小提琴、中提琴、大提琴、低音提琴组成。它们在音区、音色和角色上不同，不能把五条轨都写成“同样的和弦”：

\begin{itemize}
\item \textbf{第一小提琴}：音区最高，常承担主旋律或最高声部；旋律清楚时，它是听众先听到的那条线。
\item \textbf{第二小提琴}：低于第一小提琴，常填内声部、加倍主旋律的低八度，或与中提琴组成对答。
\item \textbf{中提琴}：中频的“胶水”，把高音小提琴和低音大提琴接起来；它的音色比小提琴暗，适合内声部和声。
\item \textbf{大提琴}：既能守低音，也能唱旋律；在中音区它常是“第二支旋律”。
\item \textbf{低音提琴}：把根音钉在最低处，让整个弦乐组有地板；它不适合承担快速旋律。
\end{itemize}

写弦乐配器，第一步不是选音色，而是写一张五行的\textbf{角色分配表}：哪条声部唱主旋律，哪条守低音，哪些填空。没有这张表，五条轨很容易写成五份互相打架的钢琴谱。

\section{音区：先保证每条声部都在自己的“舒服区”}
\label{sec:string-ranges}

下表是常用音区，按实际可演奏的保守范围写，不作为绝对限制：

\begin{longtable}{p{0.18\textwidth}p{0.20\textwidth}p{0.42\textwidth}}
\toprule
声部 & 常用音区 & 配器用途 \\
\midrule
\endfirsthead
\toprule
声部 & 常用音区 & 配器用途 \\
\midrule
\endhead
第一小提琴 & G3--E7 & 主旋律、最高声部；低于 G3 会失去亮度 \\
第二小提琴 & G3--C7 & 内声部、主旋律低八度、与中提琴对答 \\
中提琴 & C3--A6 & 中频连接、内声部和声；不要长期挤在最低处 \\
大提琴 & C2--A5 & 低音、次旋律、副歌时的对位 \\
低音提琴 & E1--G4 & 根音与低音锚；高音区只做短句点缀 \\
\bottomrule
\end{longtable}

\begin{figure}[htbp]
  \centering
  \begin{tikzpicture}[font=\small]
    \draw[->] (0,0) -- (9.6,0) node[right] {pitch};
    \foreach \x/\lab in {0/E1,1.2/C2,2.4/G3,3.6/C3,4.8/C4,6.0/C5,7.2/C6,8.4/C7,9.4/E7} {
      \draw[gray!40] (\x,0.1) -- (\x,3.6);
      \node[font=\scriptsize, rotate=60, anchor=east] at (\x,-0.12) {\lab};
    }
    \fill[yebanblue!70] (0.0,2.9) rectangle (8.8,3.25);
    \fill[yebanblue!45] (2.4,2.25) rectangle (8.2,2.55);
    \fill[yebanblue!30] (3.2,1.6) rectangle (7.6,1.9);
    \fill[yebanblue!55] (1.2,0.95) rectangle (7.0,1.25);
    \fill[gray!45] (0.0,0.3) rectangle (3.5,0.6);
    \node[font=\scriptsize, anchor=west] at (9.0,3.05) {Violin I};
    \node[font=\scriptsize, anchor=west] at (9.0,2.38) {Violin II};
    \node[font=\scriptsize, anchor=west] at (9.0,1.73) {Viola};
    \node[font=\scriptsize, anchor=west] at (9.0,1.08) {Cello};
    \node[font=\scriptsize, anchor=west] at (9.0,0.43) {Bass};
  \end{tikzpicture}
  \caption{弦乐组五声部的常用音区；条形只表示“能站住的范围”，不代表每首曲子都要占满。}
  \label{fig:string-ranges}
\end{figure}

音区不是越宽越好。第一小提琴写到 G3 以下，音色会闷；低音提琴写到 G4 以上，短句可以，长句会失去支撑。先把每条声部放回舒服区，再谈和声。

\section{配器骨架：主旋律、低音、内声部}
\label{sec:string-skeleton}

弦乐配器的第一层是三件事：

\begin{enumerate}
\item \textbf{主旋律}：通常交给第一小提琴，或在副歌时交给大提琴换色；不要同时让五条声部都唱同一句旋律。
\item \textbf{低音锚}：大提琴与低音提琴守住根音，低音提琴只在需要更厚时加入；两件乐器同音重复会显得笨重，低音提琴可低八度或只落强拍。
\item \textbf{内声部}：第二小提琴与中提琴填和弦的中间音，让主旋律和低音之间不断裂。
\end{enumerate}

一个八小节副歌可以这样分：

\begin{longtable}{p{0.22\textwidth}p{0.34\textwidth}p{0.34\textwidth}}
\toprule
声部 & 副歌第 1--4 小节 & 副歌第 5--8 小节 \\
\midrule
\endfirsthead
\toprule
声部 & 副歌第 1--4 小节 & 副歌第 5--8 小节 \\
\midrule
\endhead
第一小提琴 & 主旋律 & 主旋律，尾句加高八度 \\
第二小提琴 & 主旋律下方三度/六度 & 与中提琴交换内声部 \\
中提琴 & 和弦三音/五音 & 对答短句 \\
大提琴 & 根音与小节线落点 & 根音，尾句上唱一次 \\
低音提琴 & 根音，强拍 & 根音，强拍；不在弱拍抢戏 \\
\bottomrule
\end{longtable}

这张表比“和弦对不对”更重要：它决定了听众听到的是一条旋律加一片伴奏，还是五条互相遮住的线。

\section{排列：低音要宽，高音要近}
\label{sec:string-spacing}

弦乐的\term{排列}{voicing}与键盘不同：五条声部同时发声时，低音区如果挤得太近，会糊成一片；高音区如果拉得太开，会显得空。

\begin{itemize}
\item \textbf{低音区宽}：大提琴与低音提琴之间留出五度以上的空间，让根音和上方声部不打架。
\item \textbf{高音区近}：第一、第二小提琴与中提琴的音程控制在三度、六度或四度内，避免大跳造成空洞。
\item \textbf{内声部让路}：内声部不要和主旋律同音、同节奏；主旋律休止时，内声部才填进来。
\item \textbf{避免声部交叉}：第一小提琴不要低于第二小提琴，中提琴不要长期低于大提琴；交叉会让听众分不清谁是谁。
\item \textbf{声部数量}：三声部（主旋律、内声部、低音）已经够用；五声部全开只在副歌或高潮，主歌保持稀疏。
\end{itemize}

\term{分声部}{divisi}是弦乐特有的写法：一个声部内部再分成两半。它能让和声更丰富，但也会让每个声部变薄；没有把握时，先用三声部写作，再在副歌局部展开。

\section{弓法与织体：让弦乐“会说话”}
\label{sec:string-bowing}

同样的音符，用不同\term{弓法}{bowing}演奏，性格完全不同。写作时至少区分三种：

\begin{itemize}
\item \textbf{连奏}：音符之间不断开，适合铺底、抒情旋律与副歌的长句；长连奏需要呼吸，不要把八小节写成一条永不换气的线。
\item \textbf{拨奏}：短促、颗粒清楚，适合轻快段落或给节奏组加颜色；拨奏不要写得太密，否则会变成第二套鼓。
\item \textbf{顿弓/短音}：每个音符独立起落，适合有推进感的段落；短音之间要留空，否则听起来像机关枪。
\end{itemize}

\term{全奏}{tutti}与\term{室内乐式}{chamber-like}是两种密度：全奏让五条声部都发声，适合副歌与高潮；室内乐式只留三声部或更少，适合主歌与桥段。两种密度交替，比“从头到尾都是厚弦乐”更有层次。

弦乐的动态不是靠推子硬拉出来的。写作时用音符长度、密度和力度泳道做出“渐强—渐弱”的形状，混音阶段只做最后修饰。

\section{在夜半里落地}
\label{sec:string-yeban}

\begin{enumerate}
\item \textbf{先建五条 MIDI 轨}：按第一小提琴、第二小提琴、中提琴、大提琴、低音提琴命名；不要把所有弦乐写进一条轨再看不出声部。
\item \textbf{装 SFZ 音色}：用「文件 → 装载 SFZ 音色库...」给每条轨装一份弦乐音色；同一条轨只放一个声部，不要把五条声部挤进同一个音色。
\item \textbf{在钢琴卷帘里按行写}：按 P / F7 打开钢琴卷帘，先写第一小提琴的主旋律，再写大提琴的低音，最后填内声部；每次只写一条，避免五条线同时改。
\item \textbf{用力度泳道做呼吸}：在卷帘下方给长连奏画渐强/渐弱，句尾留半拍；不要用整轨推子代替乐句起伏。
\item \textbf{用自动化泳道做段落变化}：在编曲区展开自动化泳道（A），让弦乐在副歌进入时抬起、在主歌退下；参数随时间写进工程，导出时一并生效。
\item \textbf{用 Solo/Mute 检查平衡}：先 Solo 第一小提琴确认旋律独立，再逐条打开内声部，确认没有抢主旋律；最后打开节奏组，确认弦乐没有被鼓盖住。
\end{enumerate}

\screenshot{arrangement-string-voicing}{五条弦乐 MIDI 轨、钢琴卷帘与自动化泳道同屏：第一小提琴主旋律、大提琴低音、内声部分层可见}

\section{常见错误与修正}
\label{sec:string-mistakes}

\begin{longtable}{p{0.34\textwidth}p{0.58\textwidth}}
\toprule
常见错误 & 修正 \\
\midrule
\endfirsthead
\toprule
常见错误 & 修正 \\
\midrule
\endhead
五条声部都写主旋律 & 只留第一小提琴唱；其他声部改为内声部、低音或对答。 \\
低音区挤成一团 & 大提琴与低音提琴拉开五度以上；低音提琴只守根音与强拍。 \\
每条轨都写完整和弦 & 弦乐是声部组，不是五架钢琴；按主旋律、内声部、低音分工。 \\
第一小提琴写到 G3 以下 & 把低音交给中提琴/大提琴；小提琴回到 G3 以上。 \\
从主歌到副歌始终全奏 & 主歌用三声部，副歌再开五声部；密度要有增有减。 \\
连奏没有呼吸 & 每四小节留一次句尾，长句用力度泳道做出起伏。 \\
内声部与主旋律同节奏 & 内声部错开半拍或改成更长的音符，主旋律休止时才填。 \\
\bottomrule
\end{longtable}

\section{练习}
\label{sec:string-exercise}

\begin{exercise}
用八小节副歌写一段五声部弦乐：
\begin{enumerate}
\item 先写第一小提琴的主旋律与大提琴的根音，只写这两条；确认旋律与低音能独立成立。
\item 加第二小提琴与中提琴，填和弦的三音/五音；内声部不要与主旋律同节奏。
\item 低音提琴只落强拍根音；再检查五条声部有没有交叉、低音区有没有挤在一起。
\item 最后在副歌第 5--8 小节开一次全奏，前四小节保持三声部；用力度泳道给长连奏留呼吸。
\end{enumerate}
写完后用 Solo 逐条听：哪条声部可以删掉而音乐不变？能删的那条，就是配器里多余的。
\end{exercise}

\section{本章要点}
\label{sec:string-takeaways}

\begin{itemize}
\item 弦乐组是五条角色的队伍，不是一件“弦乐音色”。
\item 先分角色：主旋律、低音锚、内声部；再谈和弦。
\item 每条声部放回自己的音区；低音宽、高音近，避免声部交叉。
\item 主歌用三声部，副歌再开五声部；全奏与室内乐式交替。
\item 弓法决定性格：连奏、拨奏、短音各有用途。
\item 在夜半里用五条 MIDI 轨、SFZ、力度泳道与自动化落地；混音只做最后修饰。
\end{itemize}

\section{扩展阅读}
\label{sec:06-strings-further-reading}
\begin{itemize}
  \item \href{https://en.wikipedia.org/wiki/String_section}{String section（弦乐组）}：先看五个声部的标准编制。
  \item \href{https://en.wikipedia.org/wiki/Orchestration}{Orchestration（配器）}：看声部如何在乐队里分工。
  \item \href{https://www.youtube.com/watch?v=UlJ3a7xvCYI}{弦乐声部进行教程}：专讲和弦摆位与声部进行，避免各弹各的。
  \item \href{https://www.youtube.com/watch?v=CSq3tsfYcUU}{从钢琴草稿到弦乐配器}：用免费采样也能做出的分层写法。
\end{itemize}

\begin{helpmapping}
\begin{longtable}{p{0.45\textwidth}p{0.45\textwidth}}
\toprule
小节 & 帮助锚点 \\
\midrule
弦乐组先是一支队伍 & \texttt{arrangement-string-section} \\
音区 & \texttt{arrangement-string-ranges} \\
配器骨架 & \texttt{arrangement-string-skeleton} \\
排列 & \texttt{arrangement-string-spacing} \\
弓法与织体 & \texttt{arrangement-string-bowing} \\
在夜半里落地 & \texttt{arrangement-string-yeban} \\
常见错误与修正 & \texttt{arrangement-string-mistakes} \\
练习 & \texttt{arrangement-string-exercise} \\
\bottomrule
\end{longtable}
\end{helpmapping}
